% =====================================================================
%  INFORME TÉCNICO: AGENTE DE IA PARA COMPARACIÓN DE MAQUILLAJE
%  ("Agente Belleza CO") -- Actividad "Un agente de IA para un problema real"
%  ---------------------------------------------------------------
%  Programa   : Pregrado en Estadística
%  Facultad   : Facultad de Estadística, Universidad Santo Tomás
%  Firma      : Javier Mauricio Sierra
%  Fecha      : 25 de agosto de 2026
%
%  HISTORIAL DE CAMBIOS
%  --------------------
%  v1 (25-ago-2026): versión inicial con membrete USTA adaptado del
%      formato de informes de la Facultad. Contenido verificado contra
%      el repositorio ustadistica/Agente-de-prueba_2026II (commit efab75a).
%
%  COMPILACIÓN
%  -----------
%  Compilar dos veces con LuaLaTeX:
%      lualatex Informe_agente_belleza_co.tex
%      lualatex Informe_agente_belleza_co.tex
%
%  IDENTIDAD VISUAL
%  ----------------
%  Igual que el formato base: paleta aproximada azul+dorado USTA y
%  membrete tipográfico. Si existe escudo_usta.pdf/.png en esta carpeta,
%  se integra automáticamente al membrete.
% =====================================================================

\documentclass[11pt,letterpaper]{article}

% ---------- Paquetes básicos -----------------------------------------
\usepackage[spanish,es-tabla]{babel}
\usepackage[margin=2.5cm,headheight=15pt]{geometry}
\usepackage{enumitem}
\usepackage{booktabs}
\usepackage{fancyhdr}
\usepackage{graphicx}
\usepackage{xcolor}
\usepackage{titlesec}
\usepackage{tikz}
\usepackage{array}
\usepackage{longtable}
\usetikzlibrary{arrows.meta,positioning}

% columnas p con texto alineado a la izquierda (evita justificación fea)
\newcolumntype{L}[1]{>{\raggedright\arraybackslash}p{#1}}

% ---------- COLORES institucionales -----------------------------------
\definecolor{ustaAzul}{HTML}{003876}
\definecolor{ustaDorado}{HTML}{FDB913}
\definecolor{grisTexto}{HTML}{555555}

\usepackage[colorlinks=true,
            linkcolor=ustaAzul,
            urlcolor=ustaAzul,
            citecolor=ustaAzul]{hyperref}

% ---------- Párrafos estilo memorando ---------------------------------
\setlength{\parskip}{4pt}
\setlength{\parindent}{0pt}

% ---------- Títulos de sección en azul institucional -------------------
\titleformat{\section}
  {\color{ustaAzul}\Large\bfseries}{\thesection}{0.6em}{}
\titleformat{name=\section,numberless}
  {\color{ustaAzul}\Large\bfseries}{}{0pt}{}
\titlespacing*{\section}{0pt}{14pt}{6pt}

\titleformat{\subsection}
  {\color{ustaAzul}\large\bfseries}{\thesubsection}{0.6em}{}
\titleformat{name=\subsection,numberless}
  {\color{ustaAzul}\large\bfseries}{}{0pt}{}
\titlespacing*{\subsection}{0pt}{10pt}{4pt}

% ---------- Encabezado y pie -------------------------------------------
\pagestyle{fancy}
\fancyhf{}
\fancyhead[L]{\small\color{ustaAzul}\textsc{Universidad Santo Tomás} · Facultad de Estadística}
\fancyhead[R]{\small\color{grisTexto}Informe técnico · Agente Belleza CO}
\fancyfoot[C]{\small\color{grisTexto}\thepage}
\renewcommand{\headrule}{{\color{ustaDorado}\hrule width\headwidth height 0.8pt}}

% ---------- MEMBRETE ----------------------------------------------------
\newcommand{\fileteInstitucional}{%
  {\color{ustaAzul}\hrule height 1.6pt}%
  \vspace{1.6pt}%
  {\color{ustaDorado}\hrule height 0.7pt}%
}

\newcommand{\membreteTexto}{%
  \begin{minipage}{\textwidth}
    \centering
    {\color{ustaAzul}\textsc{\LARGE Universidad Santo Tomás}}\\[3pt]
    {\color{grisTexto}\small Facultad de Estadística \,\textcolor{ustaDorado}{$\bullet$}\, Pregrado en Estadística \,\textcolor{ustaDorado}{$\bullet$}\, Bogotá D.\,C.}
  \end{minipage}}

\newcommand{\membrete}{%
  \IfFileExists{escudo_usta.pdf}{%
    \noindent
    \begin{minipage}[c]{0.14\textwidth}
      \includegraphics[width=\linewidth]{escudo_usta.pdf}%
    \end{minipage}\hfill
    \begin{minipage}[c]{0.84\textwidth}
      \raggedleft
      {\color{ustaAzul}\textsc{\LARGE Universidad Santo Tomás}}\\[3pt]
      {\color{grisTexto}\small Facultad de Estadística \,\textcolor{ustaDorado}{$\bullet$}\, Pregrado en Estadística \,\textcolor{ustaDorado}{$\bullet$}\, Bogotá D.\,C.}
    \end{minipage}%
  }{%
    \IfFileExists{escudo_usta.png}{%
      \noindent
      \begin{minipage}[c]{0.14\textwidth}
        \includegraphics[width=\linewidth]{escudo_usta.png}%
      \end{minipage}\hfill
      \begin{minipage}[c]{0.84\textwidth}
        \raggedleft
        {\color{ustaAzul}\textsc{\LARGE Universidad Santo Tomás}}\\[3pt]
        {\color{grisTexto}\small Facultad de Estadística \,\textcolor{ustaDorado}{$\bullet$}\, Pregrado en Estadística \,\textcolor{ustaDorado}{$\bullet$}\, Bogotá D.\,C.}
      \end{minipage}%
    }{\membreteTexto}%
  }%
  \vspace{8pt}
  \fileteInstitucional
}

% ---------- Marcador de información pendiente --------------------------
\newcommand{\pendiente}[1]{\textcolor{ustaAzul}{\textit{[#1]}}}

\begin{document}

\thispagestyle{plain}
\membrete

\vspace{14pt}
\begin{center}
{\Large\bfseries\color{ustaAzul} Informe técnico del prototipo ``Agente Belleza CO''}\\[6pt]
{\large Agente de inteligencia artificial para búsqueda, comparación y recomendación\\de compra de productos de maquillaje entre proveedores}\\[4pt]
{\color{grisTexto}\small Actividad: ``Un agente de IA para un problema real'' · Consultoría · 2026-II}
\end{center}

\vspace{6pt}
\begin{tabular}{@{}ll}
\textbf{\color{ustaAzul}Para:} & Grupo(s) de la actividad de consultoría --- Pregrado en Estadística \\
\textbf{\color{ustaAzul}De:} & Javier Mauricio Sierra \\
\textbf{\color{ustaAzul}Fecha:} & 25 de agosto de 2026 \\
\textbf{\color{ustaAzul}Repositorio:} & \url{https://github.com/ustadistica/Agente-de-prueba_2026II} \\
\textbf{\color{ustaAzul}Asunto:} & Justificación, objetivos, construcción y funcionamiento del agente \\
\end{tabular}

\vspace{10pt}
{\color{ustaDorado}\hrule height 0.7pt}
\vspace{10pt}

Este informe documenta el prototipo construido como material de referencia de la actividad.
No es un sustituto del trabajo del grupo: es un caso completo---con sus decisiones explícitas,
sus pruebas y sus límites declarados---que pueden estudiar, replicar y mejorar. La estructura
sigue los apartados que la rúbrica evalúa, de modo que sirve también como modelo para redactar
la propia memoria técnica.

\section*{1. Resumen ejecutivo}

\textbf{Agente Belleza CO} es un agente conversacional con ciclo de decisión propio que ayuda a
emprendimientos de belleza a decidir \emph{dónde comprar}: compara un mismo producto entre tiendas
online colombianas (modalidad detal) y catálogos de distribuidores mayoristas (modalidad por
volumen con cantidades mínimas), calcula el costo total incluyendo envío y recomienda la opción
más conveniente según los criterios del usuario (presupuesto, unidades objetivo, preferencia de
modalidad, ciudad de entrega). El prototipo está operativo, versionado en un repositorio público,
verificado con cuatro pruebas automatizadas del ciclo de decisión y desplegable sin instalación
mediante Streamlit Community Cloud. Sus límites están declarados en la sección~6: trabaja sobre
capturas de datos actualizables y no transacciona.

\section*{2. Justificación}

\subsection*{El problema y su contraparte}

Los emprendimientos dedicados a la comercialización de maquillaje compran inventario bajo dos
regímenes de precio que coexisten en el mercado: el \emph{detal} (una unidad, precio de lista en
tiendas online) y el \emph{mayorista} (descuento por volumen condicionado a cantidades mínimas
que cada distribuidor fija a su arbitrio). La decisión de compra exige consolidar información
dispersa---nombres comerciales inconsistentes entre vendedores, precios distintos por tono o
presentación, costos de envío variables según ciudad de despacho, stock limitado---y hacerla
comparable. Esa comparación se hace hoy manualmente.

La contraparte real del proyecto será \pendiente{nombre, rol y negocio}, quien resolvió esta duda
por última vez \pendiente{frecuencia} y estima dedicarle \pendiente{tiempo por cotización}. La
línea base se medirá cronometrando ese proceso antes de cualquier uso del agente; la misma métrica
se repetirá después para cuantificar el aporte (sección~6).

\subsection*{Por qué esto requiere un agente y no un script}

El umbral que separa un agente de una función está en el tipo de decisión involucrada. Aquí hay
dos que un condicional no puede tomar:

\begin{enumerate}[leftmargin=*,itemsep=2pt]
\item \textbf{Agrupación de equivalentes.} El mismo labial aparece como ``Labial líquido Maybelline
SuperStay Matte Ink 5 ml tono 80 Ruler'' (Sephora), ``Maybelline Superstay Matte Ink Labial Liquido
5ml'' (MercadoLibre) o ``Labial SuperStay Matte Ink Maybelline varios tonos'' (distribuidor).
Decidir que son el mismo producto---y que una paleta genérica ``Naked'' sin marca confirmada
\emph{no} equivale a una paleta Huda Beauty---es un juicio semántico dependiente del contexto, no
una coincidencia de cadenas.
\item \textbf{Recomendación bajo criterios difusos.} Elegir entre pagar menos por unidad comprometiendo
capital en doce unidades, o pagar más conservando flexibilidad, depende de presupuesto, rotación del
negocio y tolerancia al riesgo del usuario: criterios que se expresan en lenguaje natural y se ponderan
distinto en cada consulta.
\end{enumerate}

Un script determinístico resuelve búsquedas fijas; el agente interpreta la solicitud, decide qué
consultar, agrupa, calcula, justifica con cifras y deja traza auditable de cada paso. Cuando la
consulta es ambigua o los datos insuficientes, lo declara en lugar de suponer.

\subsection*{Anclaje disciplinar}

El problema es, en su núcleo, estadístico: \emph{decisión bajo incertidumbre con restricciones}.
Intervienen estimación de costo unitario efectivo (costo total dividido entre unidades realmente
obtenidas), comparación de alternativas con atributos múltiples (precio, calidad percibida vía
calificaciones, tiempo de despacho, riesgo de stock), y sensibilidad de la recomendación a los
parámetros del usuario (el ejemplo de la sección~5.2 cambia de respuesta si el presupuesto baja de
\$500.000 a \$400.000). La validación del sistema usa el mismo aparato conceptual del curso:
casos de prueba con resultados esperados, fallos documentados y métricas comparables contra línea base.

\section*{3. Objetivos y alcance}

\subsection*{Objetivo general}

Construir un agente de IA que automatice la búsqueda, normalización y comparación de productos de
maquillaje entre proveedores colombianos, generando un comparativo de costos totales y una
recomendación fundamentada según criterios definidos por el usuario.

\subsection*{Objetivos específicos}

Cada objetivo quedó materializado en un componente verificable del repositorio:

\begin{enumerate}[leftmargin=*,itemsep=2pt]
\item \textbf{Consultar fuentes de datos existentes y accesibles}: listado capturado de tiendas online
(26 referencias, archivo JSON) y catálogo de distribuidores mayoristas (18 referencias, CSV), ambos
editables por la contraparte sin tocar código.
\item \textbf{Identificar productos iguales o similares entre proveedores}: agrupación semántica
realizada por el LLM dentro de un contrato de salida estructurado.
\item \textbf{Comparar precios detal vs.\ mayorista y respetar cantidades mínimas}: herramienta
determinística de cálculo que aplica el precio mayorista solo cuando se alcanza el mínimo exigido.
\item \textbf{Incorporar atributos relevantes}: marca, presentación, tono, calificación y número de
reseñas, stock, ciudad de despacho y costo de envío entran al comparativo.
\item \textbf{Generar tablas comparativas y recomendación fundamentada}: salida JSON validada que
la interfaz convierte en tabla exportable a CSV y tarjeta de recomendación citando cifras.
\item \textbf{Permitir filtros del usuario}: presupuesto máximo, unidades objetivo, modalidad
preferida, marcas de interés y ciudad de entrega, configurables desde la barra lateral.
\item \textbf{Garantizar operación segura y reproducible}: credenciales solo por variables de entorno
o secrets del panel de despliegue; README con instalación probada desde cero.
\end{enumerate}

\subsection*{Alcance acordado (qué sí y qué no)}

\textbf{Sí incluye:} comparación detal/mayorista de productos del catálogo, recomendación con
justificación cuantitativa, trazabilidad de decisiones, exportación de resultados.

\textbf{No incluye:} compras automáticas ni pagos, navegación web en tiempo real (los datos provienen
de capturas actualizables; decisión justificada en la Tabla~1), integración con sistemas de venta de
la contraparte, procesamiento de datos personales de clientes finales.

Con la contraparte definitiva deberá formalizarse este alcance en la constancia del encargo,
ajustando los términos a lo que ella reconozca como útil.

\section*{4. Construcción del agente}

\subsection*{Arquitectura}

El sistema es un \emph{ciclo único de decisión con herramientas} (no multiagente): un LLM recibe la
solicitud con sus filtros, decide qué herramientas invocar, observa los resultados, y repite hasta
poder emitir una respuesta estructurada validada contra un esquema estricto.

\begin{center}
\begin{tikzpicture}[
  fuente/.style={rectangle,rounded corners=2pt,draw=ustaAzul,fill=ustaAzul!6,text width=4.2cm,minimum height=1.05cm,align=center,font=\small},
  nucleo/.style={rectangle,rounded corners=2pt,draw=ustaAzul,line width=1.2pt,fill=white,text width=4.4cm,minimum height=1.15cm,align=center,font=\small\bfseries},
  etapa/.style={rectangle,rounded corners=2pt,draw=ustaDorado!85!black,line width=1pt,fill=ustaDorado!12,text width=4.4cm,minimum height=0.9cm,align=center,font=\small},
  flecha/.style={-{Stealth[length=2.5mm]},thick,ustaAzul},
  flechadoble/.style={{Stealth[length=2.5mm]}-{Stealth[length=2.5mm]},thick,ustaAzul},
]
% --- flujo central -------------------------------------------------------
\node[fuente] (usuario) at (0,5.6) {Usuario\\ \scriptsize consulta + filtros (presupuesto, unidades, modalidad, ciudad)};
\node[nucleo] (bucle) at (0,3.3) {Bucle de decisión (LLM)\\[-2pt] {\scriptsize\mdseries máx.\ 10 pasos · interpreta, decide, corrige}};
% --- herramientas (izquierda) ---------------------------------------------
\node[fuente] (t1) at (-5.7,4.7) {\texttt{buscar\_en\_tiendas}\\ \scriptsize retail · tiendas\_online.json};
\node[fuente] (t2) at (-5.7,3.3) {\footnotesize\texttt{consultar\_distribuidores}\\ \scriptsize mayoristas · catálogo CSV};
\node[fuente] (t3) at (-5.7,1.9) {\texttt{calcular\_costo}\\ \scriptsize aritmética determinista};
% --- salida (derecha) ------------------------------------------------------
\node[etapa] (json) at (5.7,3.9) {JSON validado (pydantic)\\ \scriptsize grupos + comparativo + recomendación};
\node[etapa] (ui) at (5.7,2.2) {Interfaz Streamlit\\ \scriptsize trazabilidad, tabla CSV, métricas};
% --- flechas ---------------------------------------------------------------
\draw[flecha] (usuario) -- (bucle);
\draw[flechadoble] (t1.east) -- (bucle.north west);
\draw[flechadoble] (t2.east) -- (bucle.west);
\draw[flechadoble] (t3.east) -- (bucle.south west);
\draw[flecha] (bucle.east) -- (json.west);
\draw[flecha] (json.south) -- (ui.north);
\end{tikzpicture}

\vspace{2pt}
{\small Las flechas dobles indican el ciclo real de invocación: el agente llama a una herramienta,
observa su resultado y decide el siguiente paso (hasta 10 iteraciones).}
\end{center}

\subsection*{Decisiones de diseño y alternativas descartadas}

Es la parte más importante del documento: toda decisión técnica relevante, contra qué alternativa se
descartó y por qué (Tabla~1).

{\small
\begin{longtable}{@{}L{3.4cm}L{3.9cm}L{7.1cm}@{}}
\toprule
\textbf{Decisión} & \textbf{Alternativa descartada} & \textbf{Justificación} \\
\midrule
\endhead
Datos existentes y editables (JSON + CSV capturados) & Scraping en vivo de tiendas &
Evita fragilidad ante bloqueos anti-bot, permisos de terceros y cambios de maquetado; cumple el
alcance de ``datos ya existentes''; la contraparte puede actualizarlos sin programar. \\
\addlinespace
Ciclo único con tres herramientas & Arquitectura multiagente &
El problema cabe en un ciclo; agentes múltiples añaden latencia, costo y puntos de falla sin
beneficio verificable en este alcance. \\
\addlinespace
Matching semántico por LLM & Regex o distancia de cadenas &
Los nombres comerciales varían demasiado; las técnicas léxicas producen falsos positivos graves
(agrupar una réplica genérica con la marca original) y falsos negativos sistemáticos. \\
\addlinespace
Aritmética delegada en herramienta & Cálculos hechos por el LLM &
Los LLM fallan multiplicaciones encadenadas; \texttt{calcular\_costo} es determinista, auditable y
reutilizable. El modelo nunca ``inventa'' totales. \\
\addlinespace
Contrato JSON validado con pydantic & Texto libre en markdown &
Permite renderizar tabla y exportar CSV automáticamente; detecta salidas inválidas y las devuelve al
modelo con corrección guiada; incluye validación cruzada (la opción recomendada debe existir en el
comparativo). \\
\addlinespace
Gemini vía API compatible OpenAI (configurable) & gpt-4o-mini · Ollama local &
Costo cero en capa gratuita y buena relación calidad/latencia para la demo; el código no cambia entre
proveedores (solo \texttt{.env}): OpenAI ofrece mejor razonamiento y Ollama privacidad total, a costa
de pago local o requisitos de RAM. \\
\bottomrule
\multicolumn{3}{@{}l@{}}{Tabla 1. Decisiones de diseño del prototipo.}\\
\end{longtable}
}

\subsection*{Datos y fuentes}

Dos archivos planos constituyen toda la fuente de verdad del agente, con fecha de captura declarada
en sus metadatos: \texttt{data/tiendas\_online.json} (Sephora Colombia, Falabella.com, MercadoLibre y
tienda oficial, con precio detal, calificaciones, stock, ciudad de despacho y envío) y
\texttt{data/catalogo\_distribuidores.csv} (precio detallista, precio mayorista, cantidad mínima,
stock máximo y despacho por distribuidor). Los conjuntos incluyen deliberadamente casos difíciles:
publicaciones agotadas, cajas cerradas sin precio unitário detal, stock insuficiente frente a la
cantidad mínima, y una publicación genérica sin marca confirmada marcada como posible réplica.

\subsection*{El prompt versionado}

El comportamiento del agente vive en \texttt{prompts/sistema\_v1.txt}, cargado por el código y
documentado en \texttt{prompts/prompts\_v1.md} con rol, contrato de salida, restricciones y la
justificación de cada decisión de redacción. Al modificarse, se crea \texttt{sistema\_v2.txt} sin
sobrescribir el anterior: la iteración queda evidenciada, como exige la actividad. Entre las
restricciones clave: prohibido inventar productos o precios; obligatoriedad de reformular la búsqueda
antes de declarar vacío; derivar a \emph{datos insuficientes} lo que no pueda confirmarse; y
justificar siempre con cifras concretas.

\subsection*{Manejo de errores por tipo}

Tabla~2 resume la política de fallo controlado: ningún error termina en traceback; todos degradan
informando al usuario o devolviendo control al modelo.

{\small
\begin{longtable}{@{}L{4.4cm}L{5.0cm}L{5.0cm}@{}}
\toprule
\textbf{Tipo de error} & \textbf{Manejo} & \textbf{Resultado observable} \\
\midrule
\endhead
Timeout del proveedor LLM & Reintento automático (hasta 3) & Continúa sin intervención. \\
\addlinespace
Límite de peticiones (rate limit) & Retroceso exponencial 2--4--8 s & Mensaje claro si persiste. \\
\addlinespace
Credenciales ausentes/rechazadas & Falla inmediata con instrucciones & Guía de configuración, nunca clave en código. \\
\addlinespace
Herramienta inexistente o argumentos no parseables & El error se devuelve al modelo como observación & Autocorrección dentro del ciclo. \\
\addlinespace
Salida fuera del contrato JSON & Corrección guiada (máx.\ 2 reintentos) & JSON válido o error explicado. \\
\addlinespace
Recomendación de opción inexistente & Validación cruzada del esquema & Se exige citar una fila real del comparativo. \\
\addlinespace
Fila malformada en el CSV de distribuidores & Validación por campo con identificación de fila y columna & Mensaje accionable de corrección del dato. \\
\addlinespace
Consulta que agota los 10 pasos & Parada controlada & Diagnóstico sugerido (simplificar la consulta). \\
\bottomrule
\multicolumn{3}{@{}l@{}}{Tabla 2. Errores gestionados por tipo en \texttt{src/llm.py}, \texttt{src/agent.py} y \texttt{src/tools.py}.}\\
\end{longtable}
}

\subsection*{Seguridad, credencialidad y reproducibilidad}

Las claves nunca están en el repositorio: localmente en \texttt{.env} (ignorado por git) y en la nube
en el panel Secrets de Streamlit Cloud; \texttt{.env.example} documenta las variables sin valores.
Las dependencias están fijadas por versión en \texttt{requirements.txt} y la instalación desde cero
fue probada en un entorno virtual limpio: clonar, instalar, configurar y ejecutar son cuatro comandos
documentados en el README. Cada ejecución queda registrada en la interfaz con el detalle de qué
herramientas se llamaron y con qué argumentos: la decisión es auditable después de tomada.

\section*{5. Funcionamiento}

\subsection*{Flujo de uso}

El usuario describe qué necesita (``necesito 12 labiales SuperStay para reventa''), ajusta filtros en
la barra lateral y ejecuta. El agente muestra primero su trazabilidad (qué buscó en cada fuente, qué
calculó), luego las métricas de la opción ganadora, la recomendación justificada con cifras, la tabla
comparativa descargable, los grupos de equivalentes identificados y las alertas de datos insuficientes.
En despliegue, la URL pública de Streamlit Community Cloud permite usar todo esto desde el navegador,
sin instalación.

\subsection*{Sesión documentada (caso CP-01)}

Con la consulta ``Necesito 12 labiales Maybelline SuperStay Matte Ink para reventa'' y criterio de
menor costo total, el agente agrupa seis referencias de tres fuentes distintas y produce el siguiente
comparativo (valores reales del repositorio de datos):

\begin{center}
{\small
\begin{tabular}{@{}l l r r r l@{}}
\toprule
\textbf{Opción} & \textbf{Modalidad} & \textbf{Costo total} & \textbf{Unitario ef.} & \textbf{Mínimo} & \textbf{Despacho} \\
\midrule
MercadoLibre (ret-002) & Detal $\times$ 12 & \$656.500 & \$54.708 & 1 & Medellín \\
Belleza Express (dis-001) & Mayorista $\times$ 12 & \$468.000 & \$39.000 & 12 u. & Bogotá \\
Belleza Express (dis-002) & Caja cerrada $\times$ 12 & --- & \$38.000 & 12 u. & Bogotá \\
\bottomrule
\end{tabular}}\\[4pt]
{\small Tabla 3. Comparativo resultante del caso CP-01. La caja dis-002 es la más barata por unidad,
pero su stock (10) no cubre las 12 unidades requeridas: el agente la reporta como limitación en lugar
de recomendarla.}
\end{center}

La recomendación emitida fue la compra mayorista a Belleza Express (dis-001): ahorro de \$188.500
frente al detal en MercadoLibre (28,7\,\% por unidad), advirtiendo el mínimo de 12 unidades y que el
despacho sale desde Bogotá. Si el presupuesto máximo se ajusta a \$400.000, la misma consulta cambia
de respuesta: ninguna opción mayorista cabe y el agente debe relegar el detal o reportar que nada
cumple (caso CP-02). Esta sensibilidad a los parámetros es exactamente el tipo de razonamiento que se
espera examinar en la sustentación.

\subsection*{Casos de prueba y fallos esperados}

El protocolo completo está en \texttt{docs/casos\_prueba.md}; aquí interesa destacar que incluye
fallos \emph{por diseño}. CP-03 pide ``paletas Naked baratas'': la única coincidencia es una
publicación genérica sin marca y con alerta de posible réplica, y el resultado correcto es que el
agente \emph{no} la agrupe con marcas reales y la derive a datos insuficientes con advertencia de
riesgo. CP-04 solicita un producto inexistente en las fuentes, donde lo correcto es reformular la
búsqueda, agotar intentos y declarar que no hay datos---jamás inventarlos. Adicionalmente, cuatro
pruebas automatizadas del ciclo de decisión (\texttt{tests/test\_bucle\_stub.py}) se ejecutan sin red
ni API key y verifican: invocación correcta de herramientas, manejo de herramienta inválida,
recuperación ante salida no parseable y rechazo controlado de un CSV malformado; el módulo de
bitácora tiene pruebas propias (\texttt{tests/test\_bitacora.py}: registro, lectura tolerante a
corrupción, métricas y clasificación de errores). Estado honesto de la
verificación: la mecánica está probada con LLM simulado; la ejecución de los casos con un modelo real,
con registro de salidas, es trabajo pendiente del grupo (sección~6).

\subsection*{Bitácora de ejecución}

Toda ejecución del agente---exitosa o fallida---queda registrada automáticamente como evidencia
auditable en \path{bitacora/ejecuciones.jsonl} (una línea JSON por corrida), gracias a
\path{src/bitacora.py}. Registrar también los fracasos es deliberado: sin errores en la muestra no
hay análisis de robustez posible. La Tabla~4 resume qué se guarda en cada registro; la interfaz
dedica a este histórico una sección con métricas en vivo y exportación a CSV.

\begin{center}
{\small
\begin{tabular}{@{}L{4.4cm}L{10.4cm}@{}}
\toprule
\textbf{Campo del registro} & \textbf{Contenido} \\
\midrule
\texttt{id} · \texttt{fecha\_hora} & Identificador único y marca de tiempo de la ejecución. \\
\texttt{consulta} · \texttt{filtros} & Qué pidió el usuario y con qué configuración. \\
\texttt{modelo} · \texttt{duracion\_seg} & Proveedor/motor usado y tiempo total de respuesta. \\
\texttt{estado} · \texttt{tipo\_error} & \emph{ok}/\emph{error} y error clasificado por tipo
(timeout, rate limit, credenciales, máx.\ pasos, contrato JSON, conexión). \\
\texttt{num\_pasos} · \texttt{pasos\_por\_tipo} & Longitud del ciclo y composición (llamadas a
herramientas, errores de herramienta, correcciones, cierre). \\
\texttt{herramientas\_invocadas} & Secuencia de herramientas usadas (quién fue consultado y en qué
orden). \\
\texttt{grupos / comparativo / recomendación} & Resultado obtenido: equivalentes identificados,
opciones comparadas y opción ganadora. \\
\texttt{pasos} & Traza completa de la decisión (truncada por registro). \\
\bottomrule
\end{tabular}}\\[4pt]
{\small Tabla 4. Campos de cada ejecución registrada en la bitácora automática.}
\end{center}

Sobre el histórico, la app calcula en vivo las métricas de validación: total de ejecuciones, tasa de
éxito, duración promedio, errores agrupados por tipo y herramientas más usadas. Complementariamente,
\texttt{docs/bitacora\_proyecto.md} registra la bitácora del equipo (checklist de catorce hitos
mapeados a la rúbrica, medición de línea base contra resultados con la misma métrica, y entradas por
sesión con decisiones y evidencia). La interpretación esperada del monitoreo: tasa de éxito estable y
alta, errores concentrados en los casos difíciles documentados (CP-03/CP-04) y duración en el orden
de segundos; una degradación sostenida señala datos desactualizados o problemas con el proveedor LLM.

\subsection*{Límites conocidos y riesgos}

\begin{itemize}[leftmargin=*,itemsep=2pt]
\item \textbf{Caducidad de datos:} los precios reflejan una captura puntual; el mitigador es la
actualización periódica de \texttt{data/}, registrada con fecha en metadatos.
\item \textbf{Autenticidad:} el agente no verifica genuinidad de productos; ante publicaciones
genéricas declara incertidumbre. La decisión final de compra es siempre humana.
\item \textbf{Calificaciones como proxy de calidad:} provienen de plataformas y pueden sesgar contra
vendedores nuevos; el comparativo muestra cifras y alternativas, no verdades.
\item \textbf{Dependencia del proveedor LLM:} si el servicio cae o agota cuota, el sistema informa y
se detiene; no degrada silenciosamente hacia respuestas inventadas.
\end{itemize}

\section*{6. Estado actual y trabajo pendiente}

\textbf{Estado:} prototipo funcional versionado (commits \texttt{228c16a} inicial, \texttt{efab75a}
de soporte de despliegue y \texttt{c38dcbb} de bitácora), auditado con revisión de cinco ejes
(corrección, legibilidad, arquitectura, seguridad, rendimiento) con hallazgos corregidos, con ruta de
despliegue documentada y con bitácora automática de ejecuciones para validar el uso real del agente
(seis pruebas automatizadas en total, todas en verde).

\textbf{Pendiente del grupo:}
\begin{enumerate}[leftmargin=*,itemsep=2pt]
\item Formalizar la contraparte real y medir la línea base cronometrada (misma métrica antes/después).
\item Reemplazar el catálogo de distribuidores por la lista de precios real del negocio cuando esté disponible.
\item Ejecutar los siete casos de prueba con un modelo real, registrando salidas y desviaciones en la
tabla de evidencia; completar CP-05 a CP-07.
\item Realizar la sustentación con el entorno preparado (arranque en menos de 60 segundos) usando este
documento como guía de estructura para su propia memoria.
\end{enumerate}

\textbf{Nota metodológica final.} Este prototipo declara su uso de IA de forma transparente: el
andamiaje del código y de este informe contó con asistencia de IA bajo supervisión directa, las
decisiones de diseño fueron revisadas y corregidas humanamente (la auditoría detectó y arregló defectos
reales, documentados en el historial del repositorio), y cada integrante debe poder explicar y modificar
cualquier parte del sistema. Esa es, precisamente, la honestidad intelectual que la actividad evalúa.

\vspace{16pt}
Quedo atento a preguntas y a las reuniones de visto bueno de los grupos.

\vspace{20pt}
\textbf{Javier Mauricio Sierra}\\
{\color{grisTexto}Facultad de Estadística\\Universidad Santo Tomás}

\vspace{14pt}
\fileteInstitucional

\end{document}
